\documentclass[10pt,a4paper]{report}
\usepackage[utf8]{inputenc}
\usepackage[francais]{babel}
\usepackage[T1]{fontenc}
\usepackage{amsmath}
\usepackage{amsfonts}
\usepackage{amssymb}
\usepackage{listings} % code highlights
\usepackage{fancyhdr} % headers and footers
\usepackage{graphicx}
\usepackage[left=2cm,right=2cm,top=2cm,bottom=2cm]{geometry}


\author{Anuta Christian,Givron Azim}
\title{The frisbee's aerodynamism}


% Hide chapters numbers 
\renewcommand{\thesection}{\arabic{section}}


% Set headers and footers
\pagestyle{fancy}
\fancyhf{}
\lfoot{MECA-H3001: fluid mechanics and transfer process}
\cfoot{\thepage}
\renewcommand{\headrulewidth}{0pt}   % head horizontal rule 
\renewcommand{\footrulewidth}{0.5pt} % foot horizontal rule



\begin{document}



\begin{titlepage}

\includegraphics[scale=0.5]{logo-polytech-ULB-FR.jpg}

\center 
\vspace{5cm}
\textsc{\large MECA-H3001} \\[0.5cm]
\textsc{\LARGE Fluid mechanics and transfer process} \\[1.5cm]
\textsc{\Large English report} %\\[1.5cm]

\rule{\textwidth}{1pt}

\vspace{2cm}

\textsc{\large Anuta Christian, Givron Azim}

\end{titlepage}



\tableofcontents
\newpage 
\section{Introduction}
The frisbee is one of the objects that was invented decades ago and which is still used today. It is a source of amusement for kids and grown-ups, but also lead to new sports such as the ultimate frisbee or the disc golf. The particular shape of the frisbee allows it to glide for really long distances but to make the most of this shape, the toss has to be well oriented. In order to find the optimal angle, we need to use two basic’s physical concepts. The first one is the gyroscopic effect which tend towards to stabilize the flight and is simply due to the spinning momentum transmitted through the frisbee’s toss. The second one is the aerodynamic lift which allows the frisbee to fly and is caused by its movement in the air.
\end{document}